% Co-governed and enforced under the Sovereign Integrity Protocol License (SIP License v1.1):
% https://github.com/tensorrent/ACS-Framework/blob/main/LICENSE
\documentclass[11pt,a4paper]{article}

%% ═══════════════════════════════════════════════════════════════════════════
%% VISUAL STYLE — ACS Trilogy, presentation-grade typography
%% ═══════════════════════════════════════════════════════════════════════════
\usepackage[T1]{fontenc}
\usepackage{amsmath,amssymb,amsthm}
\usepackage{geometry}
\usepackage{xcolor}
\usepackage{graphicx}
\graphicspath{{figures/}}
\usepackage{enumitem}
\usepackage{booktabs}
\usepackage[expansion=false]{microtype}
\usepackage[numbers,sort&compress]{natbib}
\usepackage{mathtools}
\usepackage{fancyhdr}
\usepackage{titlesec}
\usepackage{needspace}
\usepackage{placeins}
\usepackage{etoolbox}
\usepackage{tocloft}

\definecolor{acsprimary}{HTML}{1a237e}
\definecolor{acssecondary}{HTML}{3949ab}
\definecolor{acsaccent}{HTML}{c62828}
\definecolor{acsgold}{HTML}{b8860b}
\definecolor{acsgreen}{HTML}{2e7d32}
\definecolor{acsslate}{HTML}{37474f}
\definecolor{acslight}{HTML}{e8eaf6}
\definecolor{acsrule}{HTML}{9fa8da}

\usepackage[labelfont={bf,small,color=acsprimary!90!black},
            textfont={small,it}]{caption}
\usepackage{hyperref}
\usepackage{bookmark}

\geometry{top=1.15in,bottom=1.15in,left=1.2in,right=1.2in,
          headheight=16pt,headsep=0.25in,footskip=0.45in}

\hypersetup{colorlinks=true,linkcolor=acsprimary,citecolor=acssecondary,
            urlcolor=acsaccent,pdfpagemode=UseOutlines,bookmarksopen=true}

\pagestyle{fancy}
\fancyhf{}
\renewcommand{\headrulewidth}{0.4pt}
\renewcommand{\footrulewidth}{0pt}
\fancyheadoffset{0pt}
\fancyhead[L]{\small\color{acsslate}\itshape\leftmark}
\fancyhead[R]{\small\color{acsslate}\textit{The Inversion Arc}}
\fancyfoot[C]{\small\color{acsslate}\thepage}
\renewcommand{\sectionmark}[1]{\markboth{\S\thesection.\ #1}{}}

\titleformat{\section}{\normalfont\Large\bfseries\color{acsprimary}}
  {\thesection}{0.7em}{}[\vspace{-2pt}{\color{acsrule}\rule{\textwidth}{0.4pt}}]
\titleformat{\subsection}{\normalfont\large\bfseries\color{acssecondary}}
  {\thesubsection}{0.7em}{}
\titleformat{\subsubsection}{\normalfont\normalsize\bfseries\color{acsslate}}
  {\thesubsubsection}{0.7em}{}
\titlespacing*{\section}{0pt}{2.4ex plus 0.9ex minus 0.4ex}{1.4ex plus 0.3ex}
\titlespacing*{\subsection}{0pt}{2.0ex plus 0.7ex minus 0.3ex}{0.9ex plus 0.2ex}
\titlespacing*{\subsubsection}{0pt}{1.6ex plus 0.5ex minus 0.2ex}{0.7ex}

\theoremstyle{plain}
\newtheorem{theorem}{Theorem}[section]
\newtheorem{lemma}[theorem]{Lemma}
\newtheorem{corollary}[theorem]{Corollary}
\newtheorem{proposition}[theorem]{Proposition}
\theoremstyle{definition}
\newtheorem{definition}[theorem]{Definition}
\theoremstyle{remark}
\newtheorem*{theorem*}{Theorem}
\newtheorem{remark}[theorem]{Remark}
\newtheorem{conjecture}[theorem]{Conjecture}

\setlist[itemize]{leftmargin=1.3em,itemsep=2pt,topsep=3pt,parsep=0pt}
\setlist[enumerate]{leftmargin=1.5em,itemsep=2pt,topsep=3pt,parsep=0pt}
\setcounter{tocdepth}{2}
\setlength{\cftbeforesecskip}{2pt}
\setlength{\cftbeforesubsecskip}{0.5pt}
\widowpenalty=10000
\clubpenalty=10000
\displaywidowpenalty=10000
%

\newcommand{\DI}{\Delta\mathcal{I}}
\newcommand{\TE}{\mathrm{TE}}
\newcommand{\Order}[1]{\mathcal{O}(#1)}
\newcommand{\eps}{\varepsilon}
\newcommand{\FormF}{\mathbf{F}}
\newcommand{\FuncF}{\boldsymbol{\Phi}}
\renewcommand{\Form}{\mathbf{e}}
\providecommand{\Func}{\boldsymbol{\omega}}
\newcommand{\Lie}[2]{[#1,\,#2]}

\title{\textbf{The Inversion Arc: Holographic Resolution in Asymmetric Coupled Systems}\\[0.4em]
       \large Constraint--Attractor Dynamics, Tensegrity, and the Banach--Tarski Connection}
\author{Bradley Wallace\\[0.3em]
        \small \textit{Independent Researcher}\\[0.2em]
        \small \texttt{The Inversion Arc}}
\date{March 2026}

\begin{document}

\begin{titlepage}
\thispagestyle{empty}

\noindent{\color{acsrule}\rule{\textwidth}{0.4pt}}
\vspace{6pt}
\noindent{\footnotesize\color{acsslate}\hfill\textit{The Inversion Arc} \quad|\quad \textit{Independent Research} \quad|\quad March 2026}
\vspace{6pt}
\noindent{\color{acsrule}\rule{\textwidth}{0.4pt}}

\vspace{0.6in}

\begin{center}
  {\Huge\bfseries\color{acsprimary} The Inversion Arc}\\[0.4em]
  {\Large\color{acsslate} Holographic Resolution in Asymmetric Coupled Systems}\\[1.2em]
  {\normalsize\itshape\color{acsslate}
    Constraint\textendash Attractor Dynamics, Tensegrity, and the Banach\textendash Tarski Connection}\\[2.0em]
  {\large\bfseries Bradley Wallace}\\[0.25em]
  {\small\itshape\color{acsslate} Independent Researcher}
\end{center}

\vspace{0.5in}

\begin{center}
\begin{minipage}{0.88\textwidth}
\begin{center}
  {\bfseries\color{acsprimary}\large Abstract}\\[0.3em]
  {\color{acsrule}\rule{1.2in}{0.3pt}}
\end{center}
\vspace{0.5em}
\small\noindent
We develop the conceptual and structural consequences of the Asymmetric
Codependent System (ACS) framework introduced in~\cite{Wallace2026a}.
The \emph{inversion arc} describes how a system that solves a constraint
becomes the new constraint: the information asymmetry $\DI$ flips sign as
the system transitions from solution to dominant structure.  We propose
this as the \emph{holographic resolution principle} and argue that it is
instantiated by the Zamolodchikov $c$-theorem (2D QFT) and the
Komargodski--Schwimmer $a$-theorem (4D QFT), which serve as proved
examples of the inversion arc in quantum field theory.  The
\emph{constraint--attractor cycle} connects the inversion arc to the
gauge-algebraic results of~\cite{Wallace2026a}: the Palatini torsion-free
condition $T^a = 0$ is both the attractor and the next constraint.
Tensegrity structures provide a discrete geometric realisation of the
ACS, where zero modes of the rigidity matrix correspond to gauge freedoms.
The Banach--Tarski paradox is reinterpreted as a geometric ACS: the
non-commuting rotations generate an exponential branching of states, and
the emergent volume measures the 3rd-order holonomy.
\end{minipage}
\end{center}

\vfill

\noindent{\color{acsrule}\rule{\textwidth}{0.4pt}}
\end{titlepage}

\clearpage
\pagenumbering{roman}
\setcounter{page}{1}
\markboth{Contents}{}
\tableofcontents
\clearpage
\pagenumbering{arabic}
\setcounter{page}{1}

\section{Introduction}
\label{sec:intro}

\noindent\textit{This paper is the third in a trilogy.
Paper~A~\cite{Wallace2026a} develops the ACS framework and its
application to gauge theory (57 verification scripts).
Paper~B~\cite{Wallace2026b} applies the framework to number theory
and the Riemann Hypothesis.  This paper explores the universal
structural consequence --- the inversion arc --- and requires only
the BCH--TE morphism (Lemma~2.5 of~\cite{Wallace2026a}) and the
definitions of Form, Function, and $\DI$.}

\medskip\noindent
This paper develops the structural and conceptual consequences of the
Asymmetric Codependent System (ACS) framework, whose mathematical
foundations and primary physical application (the gauge-algebraic
selection of $\mathfrak{su}(3)$ from Palatini geometry) are presented
in~\cite{Wallace2026a}.  The number-theoretic application is developed
in~\cite{Wallace2026b}.

The core observation is that in any system where two structurally
different objects are codependent and asymmetric, the asymmetry
\emph{generates} observable pattern through the Lie bracket $[f,g]$
of the coupling operators.  Here we explore what happens when that
pattern becomes dominant: the solution becomes the constraint, and the
information asymmetry inverts.  This is the \emph{inversion arc}.

\subsection{Paper structure}
Section~\ref{sec:tensegrity} presents the tensegrity--gauge correspondence.
Section~\ref{sec:holographic} develops holographic resolution and the inversion arc.
Section~\ref{sec:constraint} presents the constraint--attractor cycle.
Section~\ref{sec:algebraic} collects the algebraic foundations: the
Killing-orthogonality theorem, the three-class spectral taxonomy of
adjoint flows, and the algebraic non-traversability theorem (the
ACS rereading of ER\,=\,EPR).
Section~\ref{sec:unified} collects the unified claims across all ACS domains.
Section~\ref{sec:BT} develops the Banach--Tarski connection.

\subsection*{Methodological note on epistemic compression}

The results in this trilogy were obtained through repeated cycles of
conjecture, explicit verification, and compression --- a discipline
formally articulated by Lakatos~\cite{Lakatos1976} as ``proofs and
refutations.''  Several previously stated claims were narrowed during
this process: a universal ``$2\pi$ inversion at three steps'' was
found to be representation-specific; the Wronskian on scalar zero
modes was found to fail the Leibniz axiom (Paper~B,
Section~5); a three-class rather than four-class spectral taxonomy
emerges from Jordan--Chevalley decomposition (Section~\ref{sec:taxonomy}).
Each correction tightens rather than weakens the framework.  We
distinguish throughout between proved theorems, conjectures explicitly
flagged as open, and disproved hypotheses retained as boundary
markers in the model space.
\section{Discrete Geometry: Tensegrity as ACS}
\label{sec:tensegrity}
%% ─────────────────────────────────────────────────────────────────────────────

\subsection{The tensegrity--gauge correspondence}

A tensegrity network has nodes $\{q_i\} \subset \mathbb{R}^n$, edges 
partitioned into cables (tension) and struts (compression), with energy:
\begin{equation}
  H = \sum_{(i,j)\in\mathrm{cables}} k_{ij}\,V_+(|q_i-q_j|, L_{ij})
    + \sum_{(i,j)\in\mathrm{struts}} k_{ij}\,V_-(|q_i-q_j|, L_{ij}),
\end{equation}
where $V_+$ penalises extension and $V_-$ penalises compression.

\begin{lemma}[Zero modes are gauge freedoms]\label{lem:zero-modes}
The zero eigenvalues of the rigidity matrix~\cite{Connelly1996} $K = R R^\top$ 
(where $R_{3i, e}$ is the unit vector of edge $e$ at node $i$) correspond to 
deformations that cost zero energy.  In the ACS correspondence, these
align with the gauge redundancies of the lattice gauge theory defined on
the same graph: directions in configuration space where $\DI = 0$.
\end{lemma}

\noindent
Numerical verification: the icosahedral tensegrity (12 bulk nodes, 6 struts) 
yields exactly 6 zero modes, matching the $\dim(SO(3)) + \dim(\mathrm{translations}) = 6$ 
gauge redundancies of a 3D lattice gauge theory.

The ACS structure in this case:
\begin{itemize}[noitemsep]
  \item Form = cable network (tension, pulls toward equilibrium: structural).
  \item Function = strut network (compression, maintains separation: dynamic).
  \item Coupling = node positions $q_i$ shared by both.
  \item Asymmetry: $V_+(r,L) \neq V_-(r,L)$ --- the cable and strut potentials 
    are structurally different nonlinear functions.
  \item Emergent pattern: the equilibrium configuration, which is 
    neither the minimiser of cables alone nor struts alone.
\end{itemize}

%% ─────────────────────────────────────────────────────────────────────────────
\section{Holographic Resolution}
\label{sec:holographic}
%% ─────────────────────────────────────────────────────────────────────────────

\subsection{The inversion arc}

\begin{definition}[Holographic Resolution]\label{def:holo}
Let $C$ be a constraint and $S$ a system that solves $C$
by reducing $\DI$.  $S$ achieves \emph{holographic resolution} of $C$ if:
\begin{itemize}[noitemsep]
  \item[(HR-1)] \textbf{Boundary completeness:}
    $I(\partial S;\, C) = H(C)$
    --- the mutual information between the boundary $\partial S$ and
    the constraint equals the full entropy of $C$;
  \item[(HR-2)] \textbf{Interior redundancy:}
    $I(\partial S;\, \mathrm{int}(S) \mid C) = 0$
    --- given $C$, the boundary carries no information about the interior
    beyond what $C$ already determines;
  \item[(HR-3)] \textbf{Outward gating:}
    $\DI(S \to \mathrm{exterior}) < 0$
    --- the boundary gates information outward, not inward.
\end{itemize}
The Ryu--Takayanagi formula $S_{\mathrm{EE}} = \mathrm{Area}(\partial A)/(4G)$
is recovered when $\DI = 0$: at information balance, the entanglement
entropy of a bulk region equals the area of its boundary in Planck units.
Condition (HR-1) is the information-theoretic restatement of the
RT formula; (HR-2) is the ``no-hair'' property of the holographic encoding;
(HR-3) is the statement that holographic screens gate access from
inside to outside, not the reverse.
\end{definition}

\begin{definition}[Self-Resolving Structure]\label{def:self-resolving}
A structure $S$ is \emph{self-resolving} if it contains an intrinsic 
measure $\mu(t)$ such that whenever $S$ begins to invert its information 
asymmetry ($\DI \to -\DI$), the value $\mu(t)$ becomes nonzero and 
is legible \emph{within} $S$'s own record --- not only detectable by 
external observers.
\end{definition}

\begin{remark}
The Bianchi identity $D_{[\mu}F_{\nu\rho]} = 0$ is the prototype 
self-resolving structure: any deviation from gauge-covariant curvature 
is immediately encoded in the field equations themselves.  
A magnetic monopole would violate $\nabla \cdot B = 0$, and this 
violation is legible in the same layer as the field --- not 
detectable only post hoc.
\end{remark}

\subsection{The inversion theorem}

\begin{theorem}[Inversion arc of access systems]\label{thm:inversion}
Let $S$ be a system solving constraint $C_0$ by achieving $\DI > 0$ 
(Form drives Function: open, outward-flowing).  As $S$ becomes the 
dominant access structure, new actors enter the field only through $S$.  
At this point $S = C_1$ is the new constraint, and:
\begin{equation}
  \DI(S\text{ as solution}) > 0 \quad\longrightarrow\quad 
  \DI(S\text{ as constraint}) < 0.
\end{equation}
The information asymmetry has inverted.  The key has become the lock.
\end{theorem}

\begin{proof}
The proof follows from the BCH--TE morphism (Lemma~2.5 of~\cite{Wallace2026a}).

\textbf{Step~1 (Role assignment).}
When $S$ is a solution, let $f = S$ (the outward-flowing operator, Form)
and $g = C_0$ (the constraining operator, Function).  By hypothesis,
$\DI(f \to g) > 0$: Form drives Function.

\textbf{Step~2 (Role reversal).}
When $S$ becomes the dominant access structure, the roles swap:
$f' = C_0'$ (the new dynamic actors, now Function) and $g' = S$
(now the constraining boundary, Form).  The operators are the same
but their positions in the ACS coupling are exchanged: $f' \leftrightarrow g'$.

\textbf{Step~3 (BCH sign reversal).}
The BCH--TE morphism gives $\DI(f, g) = \eps^2\,\mathbb{E}[f - g] + \Order{\eps^3}$.
Under role reversal $f \leftrightarrow g$:
\[
  \DI(g, f) = \eps^2\,\mathbb{E}[g - f] = -\eps^2\,\mathbb{E}[f - g]
  = -\DI(f, g) + \Order{\eps^3}.
\]
The first-order transfer entropy flips sign exactly.  Higher-order
corrections are $\Order{\eps^3}$ and do not alter the sign for
sufficiently small coupling.

\textbf{Step~4 (Conclusion).}
The sign of $\DI$ inverts when $S$ transitions from solution to
constraint.  The inversion rate $\tau^{-1}$ scales with the adoption
rate (how quickly $S$ becomes the dominant access layer).
\end{proof}

\begin{remark}[Computational verification]
Three instances of the inversion arc are verified in~\cite{Wallace2026a}:
\begin{enumerate}[noitemsep]
  \item The Ricci flow uniformisation ($41\times$ curvature reduction):
    the Ricci flow (solution) becomes the Einstein equation (constraint).
  \item The constraint hierarchy $\|G\| > \|D\| > \|H\|$ on the spin network:
    the Gauss constraint (Form) gates the diffeomorphism and Hamiltonian
    constraints (Function).
  \item The Wheeler--DeWitt attractor ($\ker(\hat{H})$ weight $50.7\%$):
    the quantum state flows to $\hat{H}\Psi = 0$, where the WdW equation
    (solution) becomes the definition of the physical Hilbert space
    (constraint).
\end{enumerate}
\end{remark}

\noindent
Historical instances (Table~\ref{tab:acs-instances}):

\begin{table}[h]
\centering
\begin{tabular}{lll}
\toprule
Domain & ACS Form & ACS Function \\
\midrule
2D QFT (c-theorem) & UV fixed point & IR RG flow \\
4D QFT (a-theorem) & UV couplings & Weyl anomaly \\
Gauge theory & vierbein $e$ & connection $\omega$ \\
Quantum gravity & triad $\tilde{E}$ & Ashtekar $A$ \\
Number theory & primes $\{p_n\}$ & Riemann zeros $\{\gamma_k\}$ \\
\bottomrule
\end{tabular}
\caption{ACS instances across domains.  The RG examples (top two rows) provide
established instances of monotone information flow (Zamolodchikov $c$-theorem~\cite{Zamolodchikov1986} and the
Komargodski--Schwimmer $a$-theorem~\cite{KomargodskiSchwimmer2011}), anchoring the framework in known results.}
\label{tab:acs-instances}
\end{table}

\subsection{The constraint--attractor cycle}
\label{sec:constraint}

The inversion arc (Theorem~\ref{thm:inversion}) has a dynamical interpretation
that unifies the attractor structure of the gauge theory, GL(4), and
quantum gravity sections of~\cite{Wallace2026a}.

In the Palatini formulation, the field equations drive the system
to $T^a = 0$ (torsion-free geometry).  This equilibrium is both
the \emph{attractor} of the ACS flow and the \emph{constraint}
on subsequent evolution: once torsion vanishes, the connection is
fully determined by the vierbein, and the system has no remaining
degrees of freedom at 1st order.

Breaking this constraint --- activating $T^a \neq 0$ --- is the
ACS analog of Form overcoming Function.  The excess torsion produces
chiral zero-modes (~\cite{Wallace2026a}), which carry the spinor
structure needed for complexification.  Those spinor degrees of
freedom self-organise into the next constraint: the colour gauge
symmetry $\mathrm{SU}(3)$, whose confinement mechanism provides
the new attractor.

The cycle is:
\[
  T^a = 0 \;\;\xrightarrow{\;\text{torsion activates}\;}\;\;
  T^a \neq 0 \;\;\xrightarrow{\;\text{chiral modes}\;}\;\;
  \text{spinor bundle} \;\;\xrightarrow{\;\text{complexify}\;}\;\;
  \mathfrak{su}(3) \;\;\xrightarrow{\;\text{confine}\;}\;\;
  \text{new constraint}
\]
Each stage is one step of the inversion arc.
The Lie bracket $[f,g]$ at each stage measures the asymmetry
driving the transition; when it vanishes, the system is at the
attractor; when it is non-zero, the system is being driven
toward the next attractor.

For quantum gravity, the same structure appears:
the Wheeler--DeWitt condition $\hat{H}|\Psi\rangle = 0$
(the WdW conjecture~\cite{Wallace2026a}) is the quantum attractor where
$\DI_Q = 0$.  Breaking this balance ($\hbar$ corrections,
$\DI_Q \neq 0$) produces the arrow of time --- the quantum
version of activating torsion.

For the Riemann Hypothesis, the attractor is the critical line
$\sigma = \tfrac{1}{2}$: the state where the prime--zero ACS is
stationary (the spectral stationarity theorem~\cite{Wallace2026b}).  Off-critical zeros
correspond to the system not having reached equilibrium --- the
variance grows (Section~\ref{sec:unified}), the attractor has
not been found.

In every domain, the constraint is the attractor, and the attractor
is the next constraint.  The ACS framework makes this cycle
computable: $\DI$ measures how far the system is from equilibrium,
and the Lie bracket $[f,g]$ measures what is driving it there.

%% ─────────────────────────────────────────────────────────────────────────────
\section{Algebraic Foundations of the Inversion Arc}
\label{sec:algebraic}
%% ─────────────────────────────────────────────────────────────────────────────

The inversion arc, as developed in Section~\ref{sec:holographic}, has
the structural form of a Lie bracket whose successive applications
saturate at finite order.  This section makes that statement precise.
We prove the Killing-orthogonality theorem (Section~\ref{sec:killing}),
which is the algebraic content of ``non-traversability'' in the ACS
framework.  We classify adjoint flows by spectral type
(Section~\ref{sec:taxonomy}).  We re-read the ER\,=\,EPR conjecture
in this language (Section~\ref{sec:erepr}).  Together these
results sharpen what the inversion arc \emph{is} mathematically,
and what it is not.

\subsection{The Killing-orthogonality theorem}
\label{sec:killing}

\begin{theorem}[Killing-orthogonality]\label{thm:killing}
For any matrix Lie algebra $\mathfrak{g}$ and any $X, Y \in \mathfrak{g}$,
\begin{equation}\label{eq:killing-orth}
  \mathrm{tr}\bigl([X, Y]\,X\bigr) = 0,
  \qquad
  \mathrm{tr}\bigl([X, Y]\,Y\bigr) = 0.
\end{equation}
Equivalently, the bracket output $B = [X, Y]$ is trace-orthogonal to
both of its inputs in the sense of the Killing form.
\end{theorem}

\begin{proof}
By trace cyclicity,
\(
  \mathrm{tr}(XYX) = \mathrm{tr}(YX^2),
\)
so
\(
  \mathrm{tr}\bigl([X,Y]\,X\bigr) = \mathrm{tr}(XYX) - \mathrm{tr}(YX^2) = 0.
\)
The second identity follows by symmetric exchange.
\end{proof}

\noindent
This identity is verified symbolically in
\texttt{src/paper\_c/killing\_orthogonality.py} (returning
$0$ in SymPy) and numerically across $1{,}000$ random pairs in each
of $\mathfrak{sl}(3,\mathbb{R})$ and $\mathfrak{sl}(4,\mathbb{R})$
with maximum residual $\sim 10^{-15}$.

\subsubsection*{Chirality hopping example}

In $\mathfrak{sl}(4,\mathbb{R})$ with the standard basis
$\{H_i, A_{ij}, S_{ij}\}$ where $H_i$ are Cartan generators,
$A_{ij}$ are antisymmetric, and $S_{ij}$ are symmetric off-diagonal,
the bracket
\begin{equation}\label{eq:chirality-hopping}
  [H_1, A_{01}] = 2\,S_{01}
\end{equation}
maps a Cartan element and an antisymmetric off-diagonal element to
a symmetric off-diagonal element --- crossing chirality sectors
exactly.  This is a numerical verification (exact integer match) of
the more general fact that the bracket of $\mathrm{sym}$ and
$\mathrm{anti}$ sectors lands in the opposite-chirality sector.
The Killing-orthogonality theorem then forces the result $S_{01}$ to
be Killing-orthogonal to both $H_1$ and $A_{01}$, which is verified
directly: $\mathrm{tr}(S_{01} H_1) = \mathrm{tr}(S_{01} A_{01}) = 0$.

\subsubsection*{The three ambiguity numbers in \texorpdfstring{$\mathfrak{sl}(4,\mathbb{R})$}{sl(4,R)}}

Given $B = [X, Y]$, the inverse problem ``recover $(X, Y)$ from $B$''
has three distinct ambiguity numbers depending on what one means by
recovery.

\begin{enumerate}[label=(\alph*)]
  \item \textbf{Passive ambiguity (subspace).}
    The Killing-orthogonality theorem forces both $X$ and $Y$ to lie
    in the codimension-1 subspace
    \(
      B^\perp = \{Z \in \mathfrak{sl}(4) : \mathrm{tr}(BZ) = 0\}.
    \)
    Its dimension is $\dim B^\perp = n^2 - 2 = 14$.
  \item \textbf{Generic active ambiguity (Jacobian kernel).}
    The bracket map $\mu : \mathfrak{sl}(4) \oplus \mathfrak{sl}(4) \to \mathfrak{sl}(4)$,
    $\mu(X, Y) = [X, Y]$ has differential
    $d\mu(\delta X, \delta Y) = [\delta X, Y] + [X, \delta Y]$.
    At a generic point $(X, Y)$, $\mathrm{rank}\,d\mu = n^2 - 1 = 15$
    and $\dim \ker\,d\mu = 15$ --- the locally surjective regime.
    Verified across $20$ random pairs in
    \texttt{src/paper\_c/orthogonal\_complement\_probe.py}.
  \item \textbf{Degenerate active ambiguity at $(H_1, A_{01})$.}
    At the specific Cartan-antisymmetric pair, where the bracket lands
    in a $1$-dimensional subspace, the differential degenerates to
    $\mathrm{rank}\,d\mu = 11$ and $\dim \ker\,d\mu = 19$.
\end{enumerate}

These three numbers measure different things and should not be
conflated.  Citing only one ``ambiguity factor'' without context
produces misleading statements.

\subsection{The three-class spectral taxonomy}
\label{sec:taxonomy}

The orbits of $\exp(t\,\mathrm{ad}_X)$ on $\mathfrak{g}$ classify by
the spectrum of $\mathrm{ad}_X$ via the Jordan--Chevalley decomposition
$X = X_s + X_n$ with $X_s$ semisimple, $X_n$ nilpotent, and
$[X_s, X_n] = 0$.  The semisimple part $X_s$ further splits over
$\mathbb{R}$ into eigenvalues that are either real-valued or
complex-conjugate pairs, yielding three distinct flow types.

\begin{theorem}[Spectral taxonomy of adjoint flows]\label{thm:taxonomy}
Any finite-dimensional real linear ACS operator $X$ admits a
Jordan--Chevalley decomposition into commuting semisimple and
nilpotent parts.  The semisimple part further partitions adjoint
flows into three classes:
\begin{enumerate}[label=(\roman*)]
  \item \textbf{Elliptic} --- $\mathrm{ad}_{X_s}$ has imaginary
    eigenvalues (paired conjugates).  $\exp(t\,\mathrm{ad}_X)$ is
    rotational and periodic.
  \item \textbf{Hyperbolic} --- $\mathrm{ad}_{X_s}$ has real eigenvalues.
    $\exp(t\,\mathrm{ad}_X)$ is exponential growth/decay; no $2\pi$ loop
    closure.
  \item \textbf{Parabolic} --- $X_s = 0$, $X_n \neq 0$.
    $\exp(t\,\mathrm{ad}_X)$ is polynomial.
\end{enumerate}
The generic case is a commuting sum of these classes acting on
distinct invariant subspaces.
\end{theorem}

\noindent
The taxonomy is mathematically standard~\cite{Jacobson1979,Bourbaki1989};
its relevance here is to clarify that the three-fold closure
$\mathrm{ad}^3 = (16/9)\,\mathrm{ad}$ in $\mathfrak{sl}(4,\mathbb{R})$
(Theorem~C of Paper~A~\cite{Wallace2026a}) is an algebraic statement
about the minimal polynomial of $\mathrm{ad}_{T_{B-L}}$ --- not a
geometric ``$2\pi$ inversion at three steps.''  The latter applies
only in the elliptic case.

\begin{table}[h]
\centering\small
\begin{tabular}{@{}llll@{}}
\toprule
ACS instance & Spectrum of $\mathrm{ad}_X$ & Class & Geometric content \\
\midrule
$\mathrm{SU}(2)$ quaternion & $\{\pm i\}$ & Elliptic & $2\pi$ rotation $= -I$ \\
Frenet--Serret operator & $\{0, \pm i\sqrt{\kappa^2{+}\chi^2}\}$ & Elliptic & curvature/torsion rotation \\
$\mathfrak{sl}(4,\mathbb{R})$ Palatini ($T_{B-L}$) & $\{0, \pm 4/3\}$ & Hyperbolic & boost-like, no $2\pi$ closure \\
Core rope ring ($R^3 = R$) & $\{-1, 0, +1\}$ & Hyperbolic & $\mathbb{Z}/3$ permutation \\
\bottomrule
\end{tabular}
\caption{Four ACS instances classified by spectral type.  Elliptic
systems exhibit literal $2\pi$ inversion under three-step
composition; hyperbolic systems exhibit only algebraic Cayley--Hamilton
saturation, with no rotational closure.  The Frenet--Serret operator
$A$ satisfies $A^3 = -(\kappa^2 + \chi^2)\,A$, an elliptic-class
saturation; the Palatini adjoint satisfies $\mathrm{ad}_{T_{B-L}}^3 =
(16/9)\,\mathrm{ad}_{T_{B-L}}$, a hyperbolic-class saturation.  Each
identity is verified in \texttt{src/paper\_c/spectral\_taxonomy.py}.}
\label{tab:spectral-classes}
\end{table}

\begin{remark}[On the universal-$2\pi$ overclaim]
An earlier draft of this work proposed that ``three-step bracket
chains universally produce $2\pi$ geometric inversions'' across all
ACS systems.  Verification against $\mathrm{ad}_{T_{B-L}}$ in
$\mathfrak{sl}(4,\mathbb{R})$ shows otherwise:
$\exp(2\pi\,\mathrm{ad}_{T_{B-L}})$ has spectral norm of order
$e^{8\pi/3} \approx 4{,}348$, not unity.  The flow does not close.
What is universal is finite-order Cayley--Hamilton saturation; the
specific geometric content (rotation versus boost versus shear)
depends on the spectral class.  The corrected statement is preserved
here as Theorem~\ref{thm:taxonomy}.
\end{remark}

\begin{remark}[Grading selection from spectral class]
The spectral taxonomy has a direct consequence for $\mathbb{Z}/2$
gradings: in the companion note~\cite{WallaceN2}, it is shown that
minimising a universal class of adjoint spectral functionals over
symmetric involutions on $\mathbb{R}^n$ selects the grading
aligned with the minority eigenvalue cluster of the Cartan
generator.  For compact algebras (elliptic class), the minimiser is
the trivial grading; for noncompact split algebras (hyperbolic
class), it is a nontrivial $(k, n{-}k)$ grading.  The three-class
taxonomy of Theorem~\ref{thm:taxonomy} thus determines not only
the dynamics (rotational, exponential, polynomial) but also the
preferred grading structure compatible with each spectral class.
(See~\cite{WallaceN2} for the Coleman--Mandula constraint on
interpreting this grading as physical spacetime signature.)
\end{remark}

\subsection{Algebraic non-traversability: the ACS rereading of ER\,=\,EPR}
\label{sec:erepr}

The Maldacena--Susskind ER\,=\,EPR conjecture~\cite{MaldacenaSusskind2013}
posits that entangled quantum systems are connected by non-traversable
Einstein--Rosen bridges.  In the conjecture's original formulation,
non-traversability is imposed as an external causality constraint:
local operators on either side commute (preserving relativistic causality
across the bulk geometry).

The ACS framework derives the operational content of non-traversability
from the bracket algebra alone.  No bulk geometry is invoked.

\begin{proposition}[Algebraic non-traversability]\label{prop:non-trav}
Let $B = [X, Y]$ be the bracket output of an ACS coupling.  The
projection $\pi_X : \mathfrak{g} \to \mathrm{span}\{X\}$ defined by
\(
  \pi_X(M) = \frac{\mathrm{tr}(MX)}{\mathrm{tr}(X^2)}\,X
\)
satisfies $\pi_X(B) = 0$ identically.  No information from $B$ can
be ``traversed'' back to $X$ through scalar inner-product reconstruction.
\end{proposition}

\begin{proof}
By Theorem~\ref{thm:killing}, $\mathrm{tr}(BX) = 0$.  Hence the
projection coefficient is zero.
\end{proof}

\noindent
This is verified at machine precision in
\texttt{src/paper\_c/er\_epr\_algebraic.py} (projection coefficient
$|c| < 10^{-16}$ for random $(X, Y)$ pairs in $\mathfrak{sl}(4,\mathbb{R})$).
Combined with the standard quantum-mechanical fact that local Alice
and Bob operators on a Bell state commute (verified in the same module),
this gives the operational content of non-traversability without
invoking bulk geometry: the ACS bracket structure is intrinsically
``opaque'' from inputs to outputs in the trace inner product.

\begin{remark}[Limits of the rereading]
This algebraic non-traversability is a statement about the bracket
algebra, not a derivation of holographic entropy bounds or the
Ryu--Takayanagi formula~\cite{RyuTakayanagi2006}.  Whether the algebraic
non-traversability of ACS bracket outputs can be quantitatively
matched to AdS/CFT entanglement-wedge depth saturation is an
open question requiring construction of an explicit correspondence
map.  We flag this as a research direction, not a derived consequence.
\end{remark}

%% ─────────────────────────────────────────────────────────────────────────────
\section{Unified Claims and Open Problems}
\label{sec:unified}
%% ─────────────────────────────────────────────────────────────────────────────

\subsection{The unified object}

All five domains (gauge theory, number theory, discrete geometry, renormalisation group, and quantum gravity) are projections of a single structure: an ACS on a space where 
Form and Function are the two independent fields, with the bracket 
$[f, g]$ generating curvature and the holonomy $[[f,g],\cdot]$ 
being the emergent pattern.

\begin{table}[h]
\centering
\begin{tabular}{lllll}
\toprule
Domain & Form $\Form$ & Function $\Func$ & Bracket $[f,g]$ & Pattern \\
\midrule
Gauge theory & $F_{\mu\nu}$ & $A_\mu$ & $[A_\mu, A_\nu]$ & YM equation \\
Gravity & $e^a{}_\mu$ & $\omega^{ab}{}_\mu$ & $R = d\omega + \omega^2$ & Einstein eqs.\ \\
Number theory & $\{p_n\}$ & $\{\gamma_k\}$ & Explicit formula & RH stationarity \\
Tensegrity & Cables & Struts & Equilibrium & Zero modes \\
RG flow & UV couplings & IR couplings & $\beta$-function & $c$-theorem \\
Quantum gravity & $\tilde{E}^a_i$ & $A^i_a$ & LQG constraints & Spin foam \\
\bottomrule
\end{tabular}
\caption{All five domains as projections of the ACS framework.}
\end{table}

\subsection{Open problems}

\begin{conjecture}[Standard Model from GL(4) fiber]\label{conj:SM}
Let $\mathcal{Y}^{14}$ be the bundle of positive-definite metrics over $M^4$, 
with fiber $\mathrm{GL}(4)/O(4)$ (dimension 10) and base $M^4$ (dimension 4).  
The Lie algebra $\mathfrak{gl}(4)$ acting on each fiber satisfies:
\[
  \mathfrak{gl}(4) \supset \mathfrak{so}(3,1) \supset 
  \mathfrak{su}(2) \oplus \mathfrak{u}(1).
\]
Conjecture: with the ACS coupling $(\FormF, \FuncF) = (e, \omega)$ on 
$\mathcal{Y}^{14}$, the asymmetry bracket $[f, g]$ generates 
$\mathfrak{su}(2) \oplus \mathfrak{u}(1)$ 
as the electroweak sub-algebra of the ACS asymmetry 
(see ~\cite{Wallace2026a} for the proof that $\mathfrak{su}(3)$ 
cannot arise from the $O(4)$ fiber alone).
\end{conjecture}

\begin{conjecture}[Geometric fermions from torsion]\label{conj:fermions}
In an ACS with $(e^a{}_\mu, \omega^{ab}{}_\mu)$ and non-zero torsion 
$T^a \neq 0$, the torsion--spin coupling produces modes with half-integer 
angular momentum from purely geometric data.  No independent spin structure 
is required; the Cartan--Kibble--Sciama mechanism realises fermions as the 
1st-order ACS coupling between Form and Function.
\end{conjecture}

\begin{theorem*}[Zamolodchikov $c$-theorem~\cite{Zamolodchikov1986}, restated in ACS language]\label{thm:c-theorem}
For any unitary 2D QFT, there exists a function $C(g)$ on coupling space such that $dC/dt \leq 0$ along RG flow, with equality if and only if the theory is at a fixed point.  At fixed points, $C$ equals the Virasoro central charge $c$.  In the ACS framework, fixed points correspond to $\DI = 0$ (information balance between UV Form and IR Function).  The 4D analogue (Komargodski--Schwimmer $a$-theorem~\cite{KomargodskiSchwimmer2011}) extends this monotonicity to four-dimensional QFT.
\end{theorem*}

\subsection{Gaps revealed by formalisation (updated status)}

Several gaps identified in the original draft have been addressed,
with varying levels of rigour.  We distinguish three epistemic tiers:
\emph{theorem} (rigorous proof from stated axioms),
\emph{derived} (conditional result with numerical verification
and mechanism identified), and \emph{conjecture} (matching pattern
without full derivation).

\begin{enumerate}[label=(\arabic*)]
  \item \textbf{The SM generation conjecture: conjecture.}
    The companion paper~\cite{Wallace2026a} identifies one full SM
    generation from the GL(4) fiber: all 8 electric charges and
    8 hypercharges match, including a right-handed neutrino.
    Three generations follow from Jacobi truncation of the BCH at order~3.
    The Koide angle $\theta_0 = \arctan(\lambda_{\mathrm{Wolfenstein}})$
    to $0.23\%$, linking the lepton mass hierarchy to the Cabibbo angle.
    \emph{Not yet: a field-theoretic derivation showing that the BCH
    orders map to distinct fermion families in the quantum theory.}

  \item \textbf{The Wheeler--DeWitt attractor: derived.}
    A 3-node spin network Lindblad steady state concentrates in
    $\ker(\hat{H})$ with $\langle\hat{H}\rangle = 3 \times 10^{-6}$.
    This is a numerical result on a toy model; the full spin-foam
    limit is open.

  \item \textbf{The Higgs potential shape: derived.}
    The $\DI$ landscape across the GL(4) fiber produces the Mexican Hat
    potential for $15.9\%$ of Palatini directions.  The quartic coupling
    $\lambda = 2\sqrt{3}/27$ matches the SM value to $0.85\%$.
    \emph{Not yet: a derivation of the electroweak scale $v$ from
    first principles.}

  \item \textbf{The neutrino see-saw: conjecture.}
    Three layers of suppression produce $m_\nu \times M_R = m_e^4/(9m_\tau^2)$
    (holds to $0.1\%$), predicting $M_R \sim 49$~keV.
    The Dodelson--Widrow mechanism overproduces this state by $\sim\!10^5$,
    requiring either subdominant DM or an alternative production mechanism.

  \item \textbf{Definition~\ref{def:holo} and ``emergent'':} formalised
    with mutual information bounds (HR-1 through HR-3).

  \item \textbf{Theorem~\ref{thm:inversion} (testability):} the inversion
    arc now has three verified physical instances (Ricci flow, spin network
    constraint hierarchy, and the Wheeler--DeWitt attractor).

  \item \textbf{Torsion coupling hierarchy: theorem.}
    All 15 generators of $\mathfrak{sl}(4,\mathbb{R})$ classified into
    Tier~0 (zero coupling, 9~generators) and Tier~2 (coupling $32/9$,
    6~generators).  The hierarchy 0:1:4 between unbroken, electroweak,
    and colour-lepton generators follows from the $T_{B{-}L}$ eigenvalue
    spectrum (Paper~A, \S 6.8).

  \item \textbf{Bosonic vacuum energy cancellation: theorem.}
    $\sum_{X \in \mathfrak{sl}(4)} \|[T_{B{-}L}, X]\|^2 \cdot K(X,X) = 0$
    exactly, in rational arithmetic.  Reduces the cosmological constant
    problem from $10^{121}$ to $\sim 10^{55}$ (Paper~A, \S 6.8).

  \item \textbf{Corrections from the exploration session (updated).}
    (a)~$\tan\beta = 1/2$ was a misidentification of the Yukawa ratio
    $\tilde{h}/h \approx 2/3$ with the VEV ratio.
    (b)~The W/Z torsion share is a conditional prediction,
    not a theorem.  (c)~$\lambda = 2\sqrt{3}/27$ is an algebraic result
    (Koide projection), not an RG fixed point.
    (d)~The $\alpha_2$ cross-coupling
    $\mathrm{Tr}(\Phi^\dagger T^A \Phi)\,\mathrm{Tr}(\Delta_R^\dagger T^A \Delta_R)$
    vanishes identically in the minimal bi-doublet because
    $T^A \Phi = 0$ for $\Phi \sim (1, 2, 2)$ (the bi-doublet is an
    $\mathrm{SU}(4)$ singlet); $\alpha_2$ is therefore forbidden by
    representation theory, not eliminated by choice.
    (e)~The $\beta_c$ term is excluded at tree level in the minimal
    model: extremization forces $\tan\beta = \pm 1$, which by the
    no-go identity
    $M_u - M_d = (h - \tilde{h})(\kappa_1 - \kappa_2)$ gives
    $V_{\mathrm{CKM}} = I$, contradicting observed quark mixing.
    Either $\beta_c = 0$ at tree level or the Higgs sector must be
    extended (Branch~B with $\Sigma \sim (15, 1, 1)$).

  \item \textbf{The irreducible boundary: 6 inputs in Branch~A
    (minimal bi-doublet).}
    The bracket algebra and consistency filters together fix the
    Pati--Salam Higgs sector to:
    \begin{itemize}[noitemsep]
      \item Locked by Palatini: $\lambda_\Phi = 2\sqrt{3}/27$,
        $2\rho_1 + \rho_2 = 16/9$, $g = 4/3$, $\tilde{h}/h = 2/3$,
        $N_{\mathrm{gen}} = 3$.
      \item Forbidden by representation theory: $\alpha_2 = 0$.
      \item Excluded at tree level: $\beta_c = 0$ (else $V_{\mathrm{CKM}} = I$).
      \item Free at tree level: $\rho_1$, $\alpha_1$.
      \item Radiatively generated: $\tan\beta$ (value pending Coleman--Weinberg).
      \item Calibrations: $v$, $v_R$.
    \end{itemize}
    Total inputs: 4 free + 2 calibrations = 6 (vs.\ $19{+}$ in the SM,
    a reduction factor of $\sim 3.2\times$).  Whether one-loop
    Coleman--Weinberg corrections fix $\tan\beta$ uniquely (reducing
    the count to 5) is open.

  \item \textbf{Three-layer self-pruning of the Higgs sector.}
    The reduction from 5 quartic couplings to 2 free + 1 radiative
    proceeds through three independent consistency filters: representation
    theory ($\alpha_2 = 0$), vacuum extremization ($\beta_c$ forces
    $\tan\beta = \pm 1$ if present), and flavor phenomenology
    ($\kappa_1 = \kappa_2 \Rightarrow M_u = M_d$).  Each operates at a
    different mathematical layer and constrains the others; this is
    not a coincidence but a structural feature of the minimal
    bi-doublet embedding.
\end{enumerate}

\section{The Inversion Arc as a Universal Principle}
\label{sec:universal}

The constraint--attractor cycle
\[
  \text{Constraint} \;\to\; \text{Solution} \;\to\;
  \text{Dominance} \;\to\; \text{Inversion} \;\to\;
  \text{New Constraint}
\]
is instantiated across every domain touched by the ACS framework.
Table~\ref{tab:inversion-universal} maps each domain to its specific
realisation of the cycle, with the computational verification
from~\cite{Wallace2026a} indicated where available.

\begin{table}[h]
\centering\small
\begin{tabular}{@{}lllll@{}}
\toprule
Domain & Constraint & Solution & Inversion & New constraint \\
\midrule
Gauge theory & $T^a = 0$ (torsion-free) & Activate torsion &
  Chiral modes & $\mathfrak{su}(3)$ confinement \\
Quantum gravity & $\hat{H}\Psi = 0$ (WdW) & $\hbar$ corrections &
  Lindblad flow & Physical Hilbert space \\
Geometry & Curved manifold & Ricci flow & $41\times$ reduction &
  Einstein equation \\
Riemann zeros & $\sigma = \tfrac{1}{2}$ (RH) & Off-critical zeros &
  $\mathcal{S}(N,X) \to \infty$ & Stationarity restored \\
Higgs sector & $V'(\phi) = 0$ & Radial perturbation &
  Sombrero trough & VEV $= v$ \\
\bottomrule
\end{tabular}
\caption{The inversion arc across domains.  Each row instantiates the
same cycle: a constraint generates a solution; the solution becomes
dominant; dominance inverts the information asymmetry; the inverted
system becomes a new constraint.  The first three rows have
computational verification in~\cite{Wallace2026a}.}
\label{tab:inversion-universal}
\end{table}

The universality of this cycle is not merely an analogy but a structural
consequence of the BCH--TE morphism: any system with codependent
Form and Function and non-zero bracket $[f, g] \neq 0$ will exhibit
sign reversal of $\DI$ under role exchange (Theorem~\ref{thm:inversion}).
The cycle terminates at $\DI = 0$ attractors: confinement in gauge
theory, $\ker(\hat{H})$ in quantum gravity, stationarity in number
theory, the VEV in the Higgs sector.

%% ─────────────────────────────────────────────────────────────────────────────
\appendix

\section{Historical Precursors: Newton, Russell, and Mersenne}
\label{sec:history}

The ACS framework has unexpected antecedents in three historical
observations, each of which maps precisely onto the mathematical structure.

\subsection{Newton's colour wheel and the su(3) singlet}

In the \textit{Opticks} (1704), Newton arranged the visible spectrum on a
circle and observed that combining three primary colours (red, green, blue)
produces white light.  In the ACS framework:
\begin{itemize}[noitemsep]
  \item The three primary colours $\leftrightarrow$ the three weights of the
    fundamental representation of $\mathfrak{su}(3)$, separated by $120^\circ$
    in the Cartan weight diagram;
  \item White $=$ R $+$ G $+$ B $\leftrightarrow$ the colour singlet
    ($\DI = 0$, confinement);
  \item Newton's colour wheel IS the $\mathfrak{su}(3)$ weight diagram,
    topologically.
\end{itemize}
The identification is not metaphorical: QCD colour charge is named after
optical colour precisely because the combinatorics match.  The ACS recovers
both from the same structure --- the Palatini bracket on GL(4).

\subsection{Russell's wavelength spiral and atomic octaves}

In 1867, Alexander Reina Russell proposed a wavelength-based spiral
arrangement of the elements, predating Mendeleev's periodic table.
Russell's key observation: chemical periodicity has an \emph{octave}
structure, with noble gases as the ``silent'' nodes.

In the ACS framework:
\begin{itemize}[noitemsep]
  \item Shell capacities $2n^2$ (2, 8, 18, 32) arise from $\mathrm{SO}(3)$
    representation theory: $\sum_{l=0}^{n-1} (2l+1) = n^2$, doubled for spin;
  \item Noble gases ($\DI = 0$) are the ACS \emph{attractors} --- closed-shell
    configurations where the electron Form and nuclear Function are in
    perfect information balance;
  \item Chemical reactivity ($\DI \neq 0$) is the departure from the attractor;
  \item Russell's spiral IS the ACS layered resolution: each ``octave'' is
    one BCH order of the atomic Hamiltonian, with the noble gas at its
    $\DI = 0$ fixed point.
\end{itemize}

\subsection{Dimensional devolution}

The ACS framework inverts the traditional micro-to-macro construction.
The fundamental object is the full $\mathfrak{sl}(4)$ fiber (15~dimensions);
particles are \emph{projections} onto our $3{+}1$D observation space.
The compression sequence
\[
  \mathfrak{sl}(4)\,[15\mathrm{D}] \to \mathfrak{su}(3)\,[8\mathrm{D}]
  \to \mathrm{Cartan}\,[2\mathrm{D}] \to \mathrm{singlet}\,[0\mathrm{D}]
\]
has been verified computationally: the Cartan projection kills all
off-diagonal structure exactly ($\max|M_{ij}| = 0.000000$ for $i \neq j$).
A perfectly ordered geodesic on $T^{15}$ has autocorrelation $0.95$ in 15D
but only $0.23$ when projected to 3D.  What we perceive as quantum
randomness is the resolution bias of low-dimensional sampling.

%% ─────────────────────────────────────────────────────────────────────────────


\section{The Banach--Tarski Paradox as a Geometric ACS}
\label{sec:BT}

The Banach–Tarski paradox is the ultimate mathematical manifestation of a non‑commutative
system generating emergent volume.  One takes a solid ball in $\mathbb{R}^3$, partitions it
into five pieces, and reassembles them using rotations and translations to obtain two
identical copies of the original ball.  The mechanism relies entirely on the fact that
the rotation group $SO(3)$ contains a free subgroup on two generators $a$ and $b$.
These generators do not commute: $ab \neq ba$.  The set of all words in $a$ and $b$ forms
an infinite binary tree; each distinct sequence gives a distinct orientation.
The exponential growth of words ($2^n$ words of length $n$) is what allows a single piece
to be “unrolled” to cover the whole ball.

Mapping this to the ACS:
\begin{itemize}[noitemsep]
  \item \textbf{Form} and \textbf{Function} are the two rotation generators $a$ and $b$.
  \item The \textbf{2nd‑order bracket} $[a,b] \neq 0$ is exactly the non‑commutativity that
    creates the free group.
  \item The \textbf{3rd‑order holonomy} is the irreducible residue of the closed loop
    $a b a^{-1} b^{-1}$ etc. — the fact that no finite combination of generators returns
    to the identity unless it is trivial.  This is the source of the paradoxical decomposition.
  \item The \textbf{emergent volume} is the extra measure (the second ball) generated by
    the non‑commutative operations.  In ACS language, the information asymmetry $\DI$
    measures precisely this geometric defect: the failure of the flow to close under
    composition of the two generators.
\end{itemize}

Thus Banach--Tarski serves as a concrete mathematical illustration of the ACS
mechanism, where the fields are rotation operators and the measure is the Lebesgue
measure on the ball.  The paradox shows that when $[f,g] \neq 0$, the system can
create new phase-space volume, and $\DI$ provides a quantitative measure of this
non-closure across the other domains considered here.

%% ─────────────────────────────────────────────────────────────────────────────

\begin{thebibliography}{99}
\bibitem{Wallace2026a}
B.~Wallace,
\textit{Colour from Gravity: SU(3) as a Closure Attractor in the Palatini Bracket},
2026.

\bibitem{Wallace2026b}
B.~Wallace,
\textit{Spectral Susceptibility and Renormalized Stability on the Riemann Critical Line},
2026.

\bibitem{Zamolodchikov1986}
A.~B.~Zamolodchikov,
\textit{Irreversibility of the flux of the renormalization group in a 2D field theory},
JETP Letters \textbf{43}(12), 730--732 (1986).

\bibitem{KomargodskiSchwimmer2011}
Z.~Komargodski and A.~Schwimmer,
\textit{On renormalization group flows in four dimensions},
Journal of High Energy Physics \textbf{2011}(12), 099 (2011).

\bibitem{Connelly1996}
R.~Connelly and W.~Whiteley,
\textit{Second-order rigidity and prestress stability for tensegrity frameworks},
SIAM Journal on Discrete Mathematics \textbf{9}(3), 453--491 (1996).

\bibitem{Jacobson1979}
N.~Jacobson,
\textit{Lie Algebras},
Dover Publications, New York (1979).

\bibitem{Bourbaki1989}
N.~Bourbaki,
\textit{Lie Groups and Lie Algebras: Chapters 1--3},
Springer-Verlag, Berlin (1989).

\bibitem{MaldacenaSusskind2013}
J.~Maldacena and L.~Susskind,
\textit{Cool horizons for entangled black holes},
Fortschritte der Physik \textbf{61}(9), 781--811 (2013).

\bibitem{RyuTakayanagi2006}
S.~Ryu and T.~Takayanagi,
\textit{Holographic derivation of entanglement entropy from AdS/CFT},
Physical Review Letters \textbf{96}(18), 181602 (2006).

\bibitem{Lakatos1976}
I.~Lakatos,
\textit{Proofs and Refutations: The Logic of Mathematical Discovery},
Cambridge University Press, Cambridge (1976).

\bibitem{WallaceN2}
B.~Wallace,
\textit{Grading Selection from Adjoint Spectral Activity
in $\mathfrak{sl}(n,\mathbb{R})$},
2026.

\end{thebibliography}
\end{document}
